\chapterHeading{Various}

\sectionHeading{Longest Increasing Subsequence}
\begin{lstlisting}[language=C++]
int lis(const vector<int>& a) {(*\hfill*)// length, O(n log n)
    vector<int> d;
    for (int x : a) {
        auto it = lower_bound(d.begin(), d.end(), x);
        if (it == d.end()) d.push_back(x);
        else *it = x;
    }
    return (int)d.size();
}
\end{lstlisting}
\subHeading{LCS via LIS}
\begin{lstlisting}[language=C++]
// LCS -> LIS:
// Store the positions of each value of A in B.
// For each A[i], add its B positions to cs in REVERSE order
//   (this handles duplicates).
// LCS length = LIS(cs).
// If each value occurs <= K times: |cs| <= N*K
// TC: O(N*K log(N*K)); worst O(N^2 log N) if one value occurs n times.
\end{lstlisting}

\sectionHeading{Josephus problem}
\begin{lstlisting}[language=C++]
int josephus(int n, int k){ (*\hfill*) 1-based survivor, O(n)
    int res=0;
    for(int i=1; i<=n; i++) res=(res+k) % i;
    return res+1;
}
\end{lstlisting}

\sectionHeading{int128}
\begin{lstlisting}[language=C++]
__int128 read() {
    __int128 x = 0, f = 1;
    char ch = getchar();

    while (ch < '0' || ch > '9') {
        if (ch == '-') f = -1;
        ch = getchar();
    }

    while (ch >= '0' && ch <= '9') {
        x = x * 10 + (ch - '0');
        ch = getchar();
    }

    return x * f;
}

void print(__int128 x) {
    if (x < 0) {
        putchar('-');
        x = -x;
    }
    if (x > 9)
        print(x / 10);
    putchar(x % 10 + '0');
}

bool cmp(__int128 x, __int128 y) {
    return x > y;
}
\end{lstlisting}


